% year: תשע"ו
% type: מבחן
% semester: ב
% moed: ג
% source: exams/מבחנים/תשעו/מועד ג תשעו.pdf
% number: 6א
% points: 10
% categories: integral-applications

\begin{question}
שרטטו את התחום הכלוא בין העקומים $x+y=0$ ו-$y=2x-x^2$ ומצאו את שטחו.
\end{question}

\begin{hint}
מצאו את נקודות החיתוך של הישר $y=-x$ עם הפרבולה, וקבעו איזו עקומה נמצאת מעל השנייה.
\end{hint}

\begin{solution}
העקום $x+y=0$ הוא הישר $y=-x$, והעקום $y=2x-x^2$ הוא פרבולה "בוכה" (פתוחה כלפי מטה) עם קודקוד $(1,1)$ החותכת את ציר $x$ ב-$0$ וב-$2$.
\textbf{נקודות חיתוך:} $-x=2x-x^2\iff x^2-3x=0\iff x=0$ או $x=3$, כלומר הנקודות $(0,0)$ ו-$(3,-3)$.
בקטע $(0,3)$: $(2x-x^2)-(-x)=3x-x^2=x(3-x)>0$, כלומר הפרבולה מעל הישר. התחום הוא הקטע שבין הישר (מלמטה) לבין קשת הפרבולה (מלמעלה) עבור $0\le x\le3$. השטח:
\[ S=\int_0^3\big[(2x-x^2)-(-x)\big]dx=\int_0^3(3x-x^2)\,dx=\left[\frac{3x^2}{2}-\frac{x^3}{3}\right]_0^3=\frac{27}{2}-9=\boxed{\frac92}. \]
\end{solution}
